% Build from the repository root:
% latexmk -pdf -interaction=nonstopmode -halt-on-error \
%   -outdir=tmp/pdfs/qop-unitarity paper/qop-unitarity.tex
\documentclass[11pt,a4paper]{article}
\usepackage[hmargin=23mm,vmargin=22mm,headheight=14pt,headsep=7mm,footskip=11mm]{geometry}
\usepackage[T1]{fontenc}
\usepackage{amsmath,amsthm}
\usepackage{newtxtext,newtxmath}
\usepackage{microtype}
\usepackage{booktabs,tabularx,array}
\usepackage{enumitem}
\usepackage{xcolor,listings,fancyhdr}
\usepackage{hyperref}
\definecolor{ink}{HTML}{173A54}
\definecolor{codebg}{HTML}{F3F5F7}
\hypersetup{colorlinks=true,linkcolor=ink,citecolor=ink,urlcolor=ink,
  pdftitle={Soundness and Expressive Completeness of QOp},
  pdfsubject={Unitary semantics, conservative type checking, and rational angles}}
\pagestyle{fancy}
\fancyhf{}
\fancyhead[L]{\small Soundness and expressive completeness of QOp}
\fancyhead[R]{\small QPP / HQP}
\fancyfoot[C]{\thepage}
\renewcommand{\headrulewidth}{0.3pt}
\setlength{\parindent}{0pt}
\setlength{\parskip}{4pt plus 1pt}
\setlength{\emergencystretch}{2em}
\setlist{nosep,leftmargin=1.5em}
\lstset{language=Haskell,basicstyle=\small\ttfamily,
  backgroundcolor=\color{codebg},frame=single,rulecolor=\color{black!15},
  framerule=0.3pt,framesep=5pt,columns=fullflexible,
  keepspaces=true,showstringspaces=false,breaklines=true,
  aboveskip=7pt,belowskip=7pt}
\newcommand{\C}{\mathbb C}
\newcommand{\R}{\mathbb R}
\newcommand{\Q}{\mathbb Q}
\newcommand{\U}{\mathrm U}
\newcommand{\sem}[1]{\llbracket #1\rrbracket}
\newcommand{\code}[1]{\texttt{#1}}
\newcommand{\checked}[2]{#1\mathbin{:}\mathsf{QOp}(#2)}
\newcommand{\opnorm}[1]{\left\lVert #1\right\rVert_{\mathrm{op}}}
\newcommand{\rot}[2]{R_{#1}(#2)}
\DeclareMathOperator{\diag}{diag}
\newtheorem{theorem}{Theorem}
\newtheorem{lemma}[theorem]{Lemma}
\newtheorem{proposition}[theorem]{Proposition}
\newcolumntype{L}{>{\raggedright\arraybackslash}X}

\begin{document}
\begin{center}
{\LARGE\bfseries Soundness and Expressive\\[3pt] Completeness of QOp}\\[8pt]
{\large Hybrid Quantum Programming}\\[7pt]
{\small Course note \quad$\cdot$\quad QPP \quad$\cdot$\quad 18 October 2026}
\end{center}

\begin{abstract}
\noindent QOp is the language of unitary operators used in QPP. Its type
checker verifies dimensions and proves the Hermiticity conditions needed
for rotations, without constructing operator matrices. This note develops
the connection between these checks and the mathematical semantics of
quantum circuits. Accepted expressions denote unitary operators. With
real angle parameters the language can express every qubit unitary;
with its implemented rational parameters it can approximate every such
unitary arbitrarily closely. The proofs distinguish the expressive power
of the language from the ability of the checker to recognize valid
expressions.
\end{abstract}

\section{Operators, expressions, and types}

A quantum circuit on $n$ qubits acts on the state space
$\mathcal H_n=(\C^2)^{\otimes n}\cong\C^{2^n}$. Its reversible
evolution is described by a \emph{unitary} matrix $U$, satisfying
$U^\dagger U=UU^\dagger=I$, where $\dagger$ denotes conjugate transpose.
Unitarity preserves inner products and hence the norm of every state.
Write $\U(2^n)$ for the group of these matrices.

In QPP, a value of the Haskell datatype \code{QOp} is a finite syntax
tree describing an operator. It is distinct from the matrix that the
expression denotes. For example, \code{Compose X X} and \code{Id 1}
are different expressions with the same matrix denotation. Merely
constructing a \code{QOp} value does not establish its validity:
\code{Compose X (Id 2)} attempts to compose operators of different
dimensions.

The function \code{typeCheck} checks this syntax. We write
\[
  \checked a n
  \quad\text{when}\quad
  \code{typeCheck a = Right n},
\]
and use $\sem a$ for the exact matrix denotation of $a$ on $\mathcal H_n$.
Thus $n$ is the number of qubits, or \emph{width}, rather than the matrix
dimension. Throughout, widths are nonnegative and must fit in the
implementation's \code{Nat}, an alias for \code{Int}; the same condition
applies to intermediate width sums. The zero-qubit space is
$\mathcal H_0=\C$, so its operators are $1\times1$ matrices.

This note concerns the \code{QOp} fragment of QPP. Measurement and
initialization belong to the surrounding \code{Program} language and
are not unitary operations.

\subsection{Matrix semantics}

Composition follows mathematical order: \code{Compose a b} applies
$b$ first and denotes $\sem a\,\sem b$. Tensor product combines two
registers, while direct sum selects between two equal-sized blocks:
\[
  \sem{\code{Tensor}\;a\;b}=\sem a\otimes\sem b,
  \qquad
  \sem{\code{DirectSum}\;a\;b}
    =\begin{pmatrix}\sem a&0\\0&\sem b\end{pmatrix}.
\]
The direct sum has one more qubit than either block; the leading qubit
selects the first block when it is $0$ and the second when it is $1$.

Angles are measured in units of $\pi$. In particular,
\begin{equation}
  \sem{\code{Phase}\;q}=[e^{i\pi q}],
  \qquad
  \sem{\code{R}\;a\;q}=\rot A q:=e^{-i\pi qA/2},
  \quad A=\sem a.
  \label{eq:rotation}
\end{equation}
The half angle is part of the rotation convention. Both parameters
have Haskell type \code{Rational}. A phase has width zero, and
\code{Tensor (Phase q) a} denotes $e^{i\pi q}\sem a$ without adding
any qubits. Global phase is retained in the semantics: equality means
equality of matrices, rather than equality up to a scalar phase.

For the primitive gates, use the same letter in mathematical type for
the matrix and in typewriter type for its syntax; for example,
$X=\sem{\code{X}}$. The matrices are
\[
  X=\begin{pmatrix}0&1\\1&0\end{pmatrix},\quad
  Y=\begin{pmatrix}0&-i\\i&0\end{pmatrix},\quad
  Z=\begin{pmatrix}1&0\\0&-1\end{pmatrix},\quad
  H=\frac1{\sqrt2}\begin{pmatrix}1&1\\1&-1\end{pmatrix},
\]
and $\mathrm{SX}=\tfrac12\bigl((1+i)I+(1-i)X\bigr)$.
All five are unitary. The first four are also \emph{Hermitian}, meaning
$A^\dagger=A$; $\mathrm{SX}$ is not Hermitian, but satisfies
$\mathrm{SX}^2=X$.

The following table gives the constructor rules. All children must pass
the checker. In the table, $A=\sem a$ and
$B=\sem b$. A unary child has width $n$; binary children have widths
$m$ and $n$, respectively.
\begin{center}
\small
\begin{tabularx}{\linewidth}{@{}lclL@{}}
\toprule
Expression & Width & Denotation & Additional condition\\
\midrule
\code{Id n} & $n$ & $I_{2^n}$ & $n\geq0$.\\
\code{Phase q} & $0$ & $[e^{i\pi q}]$ & $q\in\Q$.\\
\code{X, Y, Z, H, SX} & $1$ & Matrices above & None.\\
\code{R a q} & $n$ & $e^{-i\pi qA/2}$ & If $q\ne0$, prove $A^\dagger=A$.\\
\code{C a} & $n+1$ & $I_{2^n}\oplus A$ & None.\\
\code{Permute ks} & $|\code{ks}|$ & Wire permutation & $\code{ks}$ permutes $0,\ldots,|\code{ks}|-1$.\\
\code{Tensor a b} & $m+n$ & $A\otimes B$ & None.\\
\code{DirectSum a b} & $n+1$ & $A\oplus B$ & $m=n$.\\
\code{Compose a b} & $n$ & $AB$ & $m=n$.\\
\code{Adjoint a} & $n$ & $A^\dagger$ & None.\\
\bottomrule
\end{tabularx}
\end{center}
For \code{Permute ks}, output wire $j$ takes input wire
$\code{ks}[j]$. The empty permutation is valid and denotes the
zero-qubit identity.

\subsection{Soundness and two notions of completeness}

Three properties describe the relationship between a language and its checker:
\begin{enumerate}
\item \emph{Soundness:} every accepted expression denotes a unitary
on its reported number of qubits.
\item \emph{Expressive completeness:} every unitary has at least one
accepted expression.
\item \emph{Completeness of the checker:} every expression with a
well-defined unitary denotation is accepted.
\end{enumerate}
The first is a correctness property of the checker. The second describes
the expressive power of the accepted language and depends on the allowed
angle parameters. The third describes how well the checker recognizes
equivalent ways of expressing operators. QPP provides soundness, but its
symbolic analysis is conservative: failure to prove a property need not
mean that the property is false.

\clearpage
\section{Soundness and Hermiticity}

Most constructor rules preserve unitarity directly. For instance, if
$A$ and $B$ are unitary and have the same dimension, then
$(AB)^\dagger AB=B^\dagger A^\dagger AB=I$.
Tensor products, direct sums, and adjoints have corresponding closure
properties. Rotation is the constructor that needs an additional premise.

\subsection{The rotation premise}

For real $q$, a Hermitian axis $A$ gives a unitary exponential:
\begin{equation}
  \rot Aq^\dagger\rot Aq
    =e^{i\pi qA/2}e^{-i\pi qA/2}=I.
  \label{eq:rotation-unitary}
\end{equation}
An arbitrary unitary axis does not suffice. For example, $XY=iZ$ is
unitary but not Hermitian, and
\[
  e^{-i\pi q(XY)/2}=e^{\pi qZ/2}
    =\diag(e^{\pi q/2},e^{-\pi q/2})
\]
is not unitary when $q\ne0$. Consequently, the checker accepts
\code{Compose X Y} as an operator but rejects it as the axis of a
nonzero rotation.

An accepted rotation axis is itself a checked operator. When the angle
is nonzero, the axis is therefore both unitary and Hermitian, so
$A^2=A^\dagger A=I$. In the exponential series, the even powers
$A^{2k}=I$ give a cosine term and the odd powers $A^{2k+1}=A$ give
a sine term. Consequently,
\begin{equation}
  \rot Aq=\cos(\pi q/2)I-i\sin(\pi q/2)A.
  \label{eq:involution-rotation}
\end{equation}
This identity also explains why rotations about checked axes can be
evaluated without a general matrix-exponential algorithm.

At $q=0$, the exponential is $I$ for any square matrix $A$. QPP therefore
omits the Hermiticity premise at zero angle, but still checks the entire
axis expression. A zero angle cannot conceal a dimension mismatch or
another invalid subterm.

\subsection{Symbolic evidence and the soundness theorem}

The checker proves Hermiticity by symbolic algebra, without expanding
the operator into a matrix. Besides the width, it may infer an exact
Pauli form $A=e^{i\pi s}P$, where $P$ is a tensor product of $I,X,Y,Z$,
or an \emph{adjoint-phase} relation
\begin{equation}
  A^\dagger=e^{i\pi p}A.
  \label{eq:adjoint-phase}
\end{equation}
The phases $s$ and $p$ are reduced modulo $2$. Since $P^\dagger=P$,
an exact Pauli form implies $p=-2s\pmod2$. A proved phase $p=0$
establishes Hermiticity; an unknown phase supplies no such proof.
For example, $H$ has adjoint phase $0$ although it is not a phase times
a Pauli string.

The soundness proof maintains two invariants together: the denotation
of each accepted subterm is unitary, and every fact inferred about that
subterm is correct. Appendix~\ref{sec:evidence} verifies the algebraic
rules used for the second invariant.

\begin{theorem}[Soundness]
Every accepted finite expression denotes a unitary on its reported
number of qubits: $\checked a n$ implies $\sem a\in\U(2^n)$.
\label{thm:soundness}
\end{theorem}
\begin{proof}
Proceed by structural induction on the expression, maintaining both
invariants above.
Identities, scalar phases, and the five primitive one-qubit gates are
unitary. A valid wire permutation is a bijection of the computational
basis, hence a unitary permutation matrix.

Products, tensor products, direct sums, and adjoints preserve
unitarity. The width restrictions ensure the stated power-of-two
dimensions. Control is the special case $I\oplus A$ of direct sum.

For \code{R a q}, the induction hypothesis validates the child's
evidence. If $q\ne0$, its accepted phase proof gives $A^\dagger=A$,
so~\eqref{eq:rotation-unitary} proves unitarity. If $q=0$, the
denotation is $I$ without a Hermiticity premise.
In each case, Lemma~\ref{lem:evidence} establishes correctness of the
parent's evidence, completing the simultaneous induction.
\end{proof}

\clearpage
\section{Expressive completeness with real parameters}

Expressive power depends on which numbers may occur as parameters.
First consider an idealized version of QOp in which \code{Phase} and
\code{R} admit real parameters. The width restrictions and the requirement
of a Hermitian axis for every nonzero rotation remain the same.
This mathematical language isolates the role of the constructors from
the representation of numbers in Haskell. The construction below uses
only the primitive axes \code{Y} and \code{Z}, whose Hermiticity is
independent of the parameter values.

\begin{theorem}[Exact expressive completeness]
With arbitrary real parameters, every $U\in\U(2^n)$ has a finite
expression satisfying the checker's restrictions and denoting $U$.
No ancillary qubits are required, and global phase is preserved.
\label{thm:complete}
\end{theorem}
\begin{proof}
For $n=0$, any unitary is a scalar $[e^{i\pi p}]$, represented by
\code{Phase p}. For one qubit, Euler decomposition expresses any
unitary as a scalar phase and three rotations:
\begin{equation}
  U=e^{i\pi p}\rot Z\alpha\rot Y\beta\rot Z\gamma,
  \qquad p,\alpha,\beta,\gamma\in\R.
  \label{eq:euler}
\end{equation}
The corresponding expression is
\begin{lstlisting}
Tensor (Phase p)
  (Compose (R Z alpha) (Compose (R Y beta) (R Z gamma)))
\end{lstlisting}
Every rotation axis is a primitive Pauli operator and therefore has
an immediate Hermiticity proof. The phase is a scalar, so the
expression has width one.

For $n\geq2$, arbitrary one-qubit unitaries and CNOT exactly generate
every $n$-qubit unitary by a finite circuit~\cite{barenco}.
The term \code{C X} supplies CNOT. Tensoring with identities embeds
one-qubit gates in any wire position. To apply CNOT to any ordered
pair of distinct wires, choose a wire permutation $P$ placing that
pair first and use
\[
  P^\dagger\bigl(\mathrm{CNOT}\otimes I_{2^{n-2}}\bigr)P.
\]
Both $P$ and $P^\dagger$ are represented by valid wire permutations.
All resulting gates have width $n$, so their composition is accepted.
Replacing each one-qubit gate by~\eqref{eq:euler} completes the
construction. Scalar phases in these decompositions are retained.
\end{proof}

This proof uses only identities, phases, rotations about $Y$ and $Z$,
\code{C X}, tensor products, composition, and wire permutations.
General rotation axes and the more sophisticated Hermiticity rules
make expressions more convenient; universality does not depend on them.
The theorem asserts existence of a finite expression, not a bound
polynomial in the number of qubits.

\subsection{Expressiveness and conservative checking}

The proof constructs an accepted expression for each operator. It does
not assert that every other expression for that operator is accepted.
This distinction permits a simple, conservative analysis to coexist
with an expressive language.

For example, let
\[
  a=\code{Compose (Tensor H H) (Tensor Y Y)}.
\]
Since $HY=-YH$, the two signs cancel in the tensor product:
\[
  (H\otimes H)(Y\otimes Y)
    =(HY)\otimes(HY)
    =(-YH)\otimes(-YH)
    =(Y\otimes Y)(H\otimes H).
\]
Thus the two factors commute. Both are Hermitian, so $\sem a$ is
Hermitian too. The current analysis does not establish this commutation
relation and rejects
\code{R a (1/3)}, although that expression denotes a unitary.
Theorem~\ref{thm:complete} nevertheless supplies an accepted
representation of the same operator in the real-parameter language.

\clearpage
\section{Rational parameters and approximation}

Return to the implemented language, where phases and angles use
Haskell's \code{Rational}. This type has arbitrary precision; the
restriction is to rational multiples of $\pi$, rather than to a
fixed number of bits.
Let
\[
  \mathcal D_n^{\Q}=\{\sem a\mid\checked a n
    \text{ in the rational-parameter language}\}
\]
be the set of matrices denoted by its accepted $n$-qubit expressions.

\begin{proposition}[Failure of exact completeness over the rationals]
For each fixed $n$, $\mathcal D_n^{\Q}$ is a proper subset of $\U(2^n)$.
\end{proposition}
\begin{proof}
There are only countably many finite expressions with integer and
rational parameters. Thus $\mathcal D_n^{\Q}$ is countable, whereas
$\U(2^n)$ is uncountable, including for $n=0$.
\end{proof}

Countability does not prevent arbitrarily accurate approximation.
To measure the error, use the operator norm
\[
  \opnorm M=\sup_{\|v\|_2=1}\|Mv\|_2.
\]
An operator-norm bound therefore controls the error on every normalized
input state, rather than on a single chosen input.

\begin{theorem}[Approximate expressive completeness]
For every $U\in\U(2^n)$ and every $\varepsilon>0$, there is an
accepted expression $a$ with rational parameters such that
\[
  \opnorm{\sem a-U}<\varepsilon.
\]
In particular, $\overline{\mathcal D_n^{\Q}}=\U(2^n)$ in operator norm.
\label{thm:density}
\end{theorem}
\begin{proof}
Use the finite circuit supplied by Theorem~\ref{thm:complete}, with
all one-qubit gates expanded as in~\eqref{eq:euler}. Its variable
parameters occur only in scalar phases and rotations about $Y$ or $Z$.
Replace them by sufficiently close rational numbers. The resulting
expression still passes the implemented checker, since none of the
required axis proofs depends on the parameter values.

For either Pauli axis $P\in\{Y,Z\}$, diagonalization with eigenvalues
$\pm1$, together with $|e^{ix}-e^{iy}|\leq|x-y|$, gives
\begin{align}
  \opnorm{\rot Pt-\rot Ps}&\leq\frac{\pi}{2}|t-s|,\label{eq:rot-bound}\\
  |e^{i\pi p}-e^{i\pi r}|&\leq\pi|p-r|.\label{eq:phase-bound}
\end{align}
Tensoring with identities and conjugating by wire permutations
preserve these norm differences. For products of unitary gates,
a telescoping expansion replaces one factor at a time. Since every
remaining factor has operator norm $1$, it gives
\begin{equation}
  \opnorm{U_m\cdots U_1-V_m\cdots V_1}
  \leq\sum_{j=1}^m\opnorm{U_j-V_j}.
  \label{eq:telescoping}
\end{equation}
There are finitely many parameters, so density of $\Q$ in $\R$
lets us make the sum smaller than $\varepsilon$.
Every approximating expression is exactly unitary by
Theorem~\ref{thm:soundness}. Since the unitary group is closed,
the asserted equality of closures follows.
\end{proof}

There are two separate errors to consider in an implementation.
Replacing real angles by rational ones selects a nearby unitary; the
resulting expression still denotes an exactly unitary matrix. A numerical
backend may then introduce additional error while evaluating that
expression. In particular, independently rounded matrix entries need
not satisfy $U^\dagger U=I$ exactly.
Type checking also places no bound on evaluation cost: an accepted
expression may exceed a backend's memory or arithmetic limits.

\clearpage
\appendix
\section{Preservation of symbolic evidence}
\label{sec:evidence}

The datatype and structural operations are defined in \code{QPP.Syntax}.
This appendix connects the symbolic analysis in \code{QPP.TypeCheck}
to the evidence invariant used in Theorem~\ref{thm:soundness}.
The helper module \code{QPP.TypeCheck.Pauli} implements arithmetic on
Pauli strings and their phases.

\begin{lemma}[Preservation of evidence]
Suppose the children are unitary with correct evidence, and the
parent's typing conditions hold. The analysis rules then infer only
correct exact Pauli forms and adjoint-phase facts for the parent.
\label{lem:evidence}
\end{lemma}
\begin{proof}
The rules fall into four groups.

\emph{Pauli forms.}
Leaf evidence follows from the primitive matrices. Exact Pauli
multiplication uses $XY=iZ$, $YZ=iX$, $ZX=iY$, their reversed
products, and $X^2=Y^2=Z^2=I$. Tensor products and adjoints preserve
Pauli form. Equal or opposite blocks satisfy
$A\oplus A=I\otimes A$ and $A\oplus(-A)=Z\otimes A$.
These identities justify the rules that retain an exact Pauli form.

\emph{Adjoint phases.}
Suppose $A^\dagger=e^{i\pi p}A$ and $B^\dagger=e^{i\pi r}B$.
Tensor products add the phases, adjoints negate them, and a direct
sum inherits a phase when both blocks have that same phase.
For composition, when $BA=e^{i\pi w}AB$,
\begin{equation}
  (AB)^\dagger=e^{i\pi(p+r)}BA
             =e^{i\pi(p+r+w)}AB.
  \label{eq:composition-phase}
\end{equation}
The checker proves this commutation relation from Pauli commutation,
identical or scalar operands, or disjoint wire supports. A support is a
set of wires outside which the operator acts as identity, apart from a
possible scalar factor. The computed support may include unnecessary
wires, but omits none that are required; thus disjoint computed supports
suffice to prove commutation.

\emph{Structural identities.}
Cancellation uses $AA^\dagger=I$, and the special square rule uses
$\mathrm{SX}^2=X$. A permutation that is its own inverse is Hermitian.
On checked terms, the structural operation \code{dagger} denotes the
adjoint, as follows by induction over its constructor equations. Matching
a term against the structural adjoint of another therefore justifies
cancellation. Unitary conjugation also preserves adjoint phase:
\[
  (UAU^\dagger)^\dagger=e^{i\pi p}UAU^\dagger.
\]
For the reflection rule, suppose $AK=-KA$ and $K$ is Hermitian. Then
$A\rot Kq A^\dagger=\rot Kq^\dagger$. Writing $B=\rot Kq$, this implies
$AB=B^\dagger A$ and $BA=AB^\dagger$, hence
\[
  (AB)^\dagger=e^{i\pi p}B^\dagger A=e^{i\pi p}AB,
  \qquad
  (BA)^\dagger=e^{i\pi p}AB^\dagger=e^{i\pi p}BA.
\]
Both products inherit $A$'s adjoint phase. At zero angle, $B=I$ and
the conclusion holds without requiring $K$ Hermitian.

\emph{Rotations.}
For a nonzero-angle rotation, the typing premise and the child's
correct evidence establish that the axis is unitary and Hermitian.
Equation~\eqref{eq:involution-rotation} therefore applies.
Even integral $q$ gives $\rot Aq=(-1)^{q/2}I$; odd integral $q$ gives
$\rot Aq=e^{-i\pi q/2}A$, with adjoint phase $1$ modulo $2$.
For a scalar axis $A=\pm I$, the result is the scalar
$e^{\mp i\pi q/2}I$ for every angle. These are exactly the scalar,
Pauli, and integral-angle phase facts used by the rotation rules.

Rotations about the same axis fuse according to
$\rot At\rot Au=\rot A{t+u}$. If the common axis lacks a
Hermiticity proof, both accepted angles are zero, so fusion still
gives the identity. A zero-angle rotation always yields the exact
identity, independently of its axis's evidence.
\end{proof}

\begin{thebibliography}{9}
\bibitem{barenco}
A.~Barenco, C.~H.~Bennett, R.~Cleve, D.~P.~DiVincenzo,
N.~Margolus, P.~Shor, T.~Sleator, J.~A.~Smolin, and H.~Weinfurter.
\newblock Elementary gates for quantum computation.
\newblock \emph{Physical Review A}, 52(5):3457--3467, 1995.
\newblock \href{https://doi.org/10.1103/PhysRevA.52.3457}{doi:10.1103/PhysRevA.52.3457}.
\newblock \href{https://arxiv.org/abs/quant-ph/9503016}{arXiv:quant-ph/9503016}.
\end{thebibliography}
\end{document}
